\documentclass[12pt,a4paper,twoside]{report}

\usepackage[T1]{fontenc}
\usepackage[utf8]{inputenc}
\usepackage[italian]{babel}
\usepackage{newtxtext}
\usepackage[a4paper,inner=30mm,outer=25mm,top=27mm,bottom=28mm,headheight=15pt]{geometry}
\usepackage{microtype}
\usepackage{fancyhdr}
\usepackage{enumitem}
\usepackage{booktabs}
\usepackage{tabularx}
\usepackage{xcolor}
\usepackage{listings}
\usepackage{titlesec}
\usepackage{hyperref}

\setlength{\parindent}{1.2em}
\setlength{\parskip}{0.1em}
\setlength{\emergencystretch}{2em}
\linespread{1.18}
\setlist{nosep,leftmargin=*}
\definecolor{acmeblue}{HTML}{173B57}
\definecolor{codebg}{HTML}{F5F7F9}
\definecolor{coderule}{HTML}{B8C4CE}
\pagestyle{fancy}
\fancyhf{}
\fancyhead[LE,RO]{\thepage}
\fancyhead[LO,RE]{\small\textsc{ACMEMobility}}
\renewcommand{\headrulewidth}{0.4pt}
\titleformat{\chapter}[display]
  {\normalfont\bfseries\color{acmeblue}}
  {\small\MakeUppercase{\chaptertitlename}\ \thechapter}
  {0.4ex}
  {\Huge}
  [\vspace{0.5ex}\titlerule]
\titlespacing*{\chapter}{0pt}{-18pt}{24pt}
\titleformat{\section}
  {\normalfont\Large\bfseries\color{acmeblue}}
  {\thesection}{0.7em}{}
\titlespacing*{\section}{0pt}{2.1ex plus 0.6ex minus 0.2ex}{0.8ex}
\hyphenation{sil-la-ba-zio-ne pa-ren-te-si}

\hypersetup{
    colorlinks=true,
    linkcolor=black,
    urlcolor=blue!50!black,
    pdftitle={ACMEMobility},
    pdfauthor={Mattia Albertini e Pierre Errico}
}

\lstset{
    basicstyle=\ttfamily\scriptsize,
    breaklines=true,
    frame=single,
    backgroundcolor=\color{codebg},
    rulecolor=\color{coderule},
    framesep=5pt,
    xleftmargin=0.5em,
    xrightmargin=0.5em,
    aboveskip=0.8em,
    belowskip=0.8em,
    columns=fullflexible,
    keepspaces=true,
    showstringspaces=false,
    literate={->}{$\rightarrow$}2
}

\newcommand{\capitoloSpecial}[1]{%
  \chapter*{#1}
  \addcontentsline{toc}{chapter}{#1}}
\newcommand{\capitoloNum}[1]{\chapter{#1}}
\newcommand{\code}[1]{\texttt{#1}}




\begin{document}
\begin{titlepage}
  \thispagestyle{empty}
  \begin{flushleft}
    {\small\bfseries\color{acmeblue}\MakeUppercase{Università di Bologna}}\\[0.35em]
    {\small Dipartimento di Informatica -- Scienza e Ingegneria}
  \end{flushleft}

  \vspace{2.2cm}
  \noindent{\color{acmeblue}\rule{\textwidth}{1.2pt}}\\[1.1cm]
  {\fontsize{34}{39}\selectfont\bfseries\color{acmeblue} ACMEMobility}\\[0.55cm]
  {\Large Architettura a microservizi per il vehicle sharing}\\[0.35cm]
  {\large Coreografia, orchestrazione BPMN e implementazione cloud-native}
  \\[1.1cm]
  \noindent{\color{acmeblue}\rule{0.34\textwidth}{0.6pt}}

  \vfill
  \begin{minipage}[t]{0.47\textwidth}
    {\small\color{acmeblue}\bfseries INSEGNAMENTO}\\[0.35em]
    Architetture Software a Microservizi\\
    Anno accademico 2025/2026
  \end{minipage}
  \hfill
  \begin{minipage}[t]{0.47\textwidth}
    {\small\color{acmeblue}\bfseries AUTORI}\\[0.35em]
    Mattia Albertini \hfill 0001235513\\
    \texttt{mattia.albertini3@studio.unibo.it}\\[0.7em]
    Pierre Errico \hfill 0001235580\\
    \texttt{pierre.errico@studio.unibo.it}
  \end{minipage}

  \vspace{1.7cm}
  {\small Bologna, \today}
\end{titlepage}
\pagenumbering{roman}
\tableofcontents
\clearpage
\pagenumbering{arabic}

\capitoloNum{Coreografia e sistema proiettato}

Durante questa prima fase abbiamo analizzato il testo e abbiamo iniziato a
suddividerlo per capire quali dovessero essere i singoli elementi della
coreografia, ottenendo i ruoli APP, ACME, BANK, FLEET e STATIONS. Successivamente
abbiamo scritto la coreografia associata al testo basandoci su questi servizi.

ACMEMobility è una piattaforma di vehicle sharing. L'utente può consultare
stazioni e veicoli, prenotare un mezzo con un anticipo massimo di 30 minuti
oppure iniziare un noleggio immediato mediante QR code, seguirne lo
spostamento e riconsegnarlo alla stazione assegnata. La cancellazione è
gratuita fino a cinque minuti prima dell'inizio; la cancellazione tardiva e
il mancato ritiro entro i cinque minuti di tolleranza producono un addebito di
10 euro. Al ritiro viene bloccata una cauzione di 10 euro. Il costo finale
dipende da durata e distanza e aumenta del 10\% se la batteria residua è
inferiore al 15\%.

Indichiamo con $m: a\rightarrow b$ l'invio del messaggio $m$ da $a$ a $b$.
La coreografia aggiornata distingue la prenotazione online (\code{ONLINE}) dal
ritiro immediato tramite QR code (\code{QR}); le risposte di errore sono
omesse per leggibilità.

\begin{lstlisting}[mathescape=true]
ONLINE :== (
    online_req: APP -> ACME;
    online_res: ACME -> APP;
    (
        pickup_req: APP -> ACME;
        pickup_res: ACME -> APP;
        unlock_req: ACME -> STATIONS;
        unlock_res: STATIONS -> ACME;
        bail_req: ACME -> BANK;
        bail_res: BANK -> ACME;
        start_req: ACME -> FLEET;
        start_res: FLEET -> ACME;
        (
            track_req: ACME -> FLEET;
            track_res: FLEET -> ACME
        )*;
        return_req: APP -> ACME;
        return_res: ACME -> APP;
        lock_req: ACME -> STATIONS;
        lock_res: STATIONS -> ACME;
        stop_req: ACME -> FLEET;
        stop_res: FLEET -> ACME;
        charge_req: ACME -> BANK;
        charge_res: BANK -> ACME;
        unbail_req: ACME -> BANK;
        unbail_res: BANK -> ACME
    )
    +
    (
        cancel_req: APP -> ACME;
        cancel_res: ACME -> APP;
        (
            bail_req: ACME -> BANK;
            bail_res: BANK -> ACME
        )
        +
        1
    )
)
QR :== (
    qr_req: APP -> ACME;
    qr_res: ACME -> APP;
    unlock_req: ACME -> STATIONS;
    unlock_res: STATIONS -> ACME;
    bail_req: ACME -> BANK;
    bail_res: BANK -> ACME;
    start_req: ACME -> FLEET;
    start_res: FLEET -> ACME;
    (
        track_req: ACME -> FLEET;
        track_res: FLEET -> ACME
    )*;
    return_req: APP -> ACME;
    return_res: ACME -> APP;
    lock_req: ACME -> STATIONS;
    lock_res: STATIONS -> ACME;
    stop_req: ACME -> FLEET;
    stop_res: FLEET -> ACME;
    charge_req: ACME -> BANK;
    charge_res: BANK -> ACME;
    unbail_req: ACME -> BANK;
    unbail_res: BANK -> ACME
)
MOBILITY :== (ONLINE + QR)*
\end{lstlisting}

La coreografia è sequentially connected in semantica sincrona: a ogni
composizione sequenziale, le interazioni ai due lati condividono almeno un
ruolo, normalmente ACME. Anche la scelta \code{ONLINE + QR} soddisfa le
condizioni di point of choice, poiché i rami iniziano entrambi con
un'interazione tra APP e ACME e coinvolgono gli stessi ruoli.

Restano alcune violazioni formali. In \code{ONLINE}, i rami di ritiro e
cancellazione non coinvolgono gli stessi ruoli; inoltre la scelta tra
\code{bail} e \code{1} coinvolge BANK in un solo ramo. Le iterazioni
\code{track*} e \code{MOBILITY*} presentano la medesima criticità nella
scelta tra ripetizione e uscita. Queste condizioni sono sufficienti, ma non
necessarie: manteniamo il modello perché i servizi assenti non hanno attività
funzionali nell'alternativa, evitando notifiche artificiali.

La proiezione è omomorfa rispetto agli operatori della coreografia: ogni
interazione diventa un output nel mittente, un input nel destinatario e
\code{1} negli altri ruoli. Indichiamo con $!m@R$ l'invio di $m$ a $R$ e con
$m@R$ la ricezione di $m$ da $R$. I rami \code{ONLINE} e \code{QR} sono
incorporati direttamente nella definizione di ciascun ruolo.

\begin{lstlisting}[mathescape=true]
APP :== ((
    !online_req@ACME;
    online_res@ACME;
    (
        !pickup_req@ACME;
        pickup_res@ACME;
        !return_req@ACME;
        return_res@ACME
    )
    +
    (
        !cancel_req@ACME;
        cancel_res@ACME
    )
)
+
(
    !qr_req@ACME;
    qr_res@ACME;
    !return_req@ACME;
    return_res@ACME
))*
ACME :== ((
    online_req@APP;
    !online_res@APP;
    (
        pickup_req@APP;
        !pickup_res@APP;
        !unlock_req@STATIONS;
        unlock_res@STATIONS;
        !bail_req@BANK;
        bail_res@BANK;
        !start_req@FLEET;
        start_res@FLEET;
        (
            !track_req@FLEET;
            track_res@FLEET
        )*;
        return_req@APP;
        !return_res@APP;
        !lock_req@STATIONS;
        lock_res@STATIONS;
        !stop_req@FLEET;
        stop_res@FLEET;
        !charge_req@BANK;
        charge_res@BANK;
        !unbail_req@BANK;
        unbail_res@BANK
    )
    +
    (
        cancel_req@APP;
        !cancel_res@APP;
        (
            !bail_req@BANK;
            bail_res@BANK
        )
        +
        1
    )
)
+
(
    qr_req@APP;
    !qr_res@APP;
    !unlock_req@STATIONS;
    unlock_res@STATIONS;
    !bail_req@BANK;
    bail_res@BANK;
    !start_req@FLEET;
    start_res@FLEET;
    (
        !track_req@FLEET;
        track_res@FLEET
    )*;
    return_req@APP;
    !return_res@APP;
    !lock_req@STATIONS;
    lock_res@STATIONS;
    !stop_req@FLEET;
    stop_res@FLEET;
    !charge_req@BANK;
    charge_res@BANK;
    !unbail_req@BANK;
    unbail_res@BANK
))*
STATIONS :== ((
    (
        unlock_req@ACME;
        !unlock_res@ACME;
        lock_req@ACME;
        !lock_res@ACME
    )
    +
    1
)
+
(
    unlock_req@ACME;
    !unlock_res@ACME;
    lock_req@ACME;
    !lock_res@ACME
))*
BANK :== ((
    (
        bail_req@ACME;
        !bail_res@ACME;
        charge_req@ACME;
        !charge_res@ACME;
        unbail_req@ACME;
        !unbail_res@ACME
    )
    +
    (
        (
            bail_req@ACME;
            !bail_res@ACME
        )
        +
        1
    )
)
+
(
    bail_req@ACME;
    !bail_res@ACME;
    charge_req@ACME;
    !charge_res@ACME;
    unbail_req@ACME;
    !unbail_res@ACME
))*
FLEET :== ((
    (
        start_req@ACME;
        !start_res@ACME;
        (
            track_req@ACME;
            !track_res@ACME
        )*;
        stop_req@ACME;
        !stop_res@ACME
    )
    +
    1
)
+
(
    start_req@ACME;
    !start_res@ACME;
    (
        track_req@ACME;
        !track_res@ACME
    )*;
    stop_req@ACME;
    !stop_res@ACME
))*
SYSTEM :==
    (APP)_APP
    |
    (ACME)_ACME
    |
    (STATIONS)_STATIONS
    |
    (BANK)_BANK
    |
    (FLEET)_FLEET
\end{lstlisting}

Le interazioni REST sono sincrone quando serve una risposta immediata; Kafka
disaccoppia produzione e lettura della telemetria; gli eventi SSE aggiornano
l'applicazione principale e i dati economici del visualizzatore. Quest'ultimo
interroga inoltre Query ogni due secondi per posizione e batteria. Camunda
conserva lo stato del processo di noleggio e correla messaggi asincroni
mediante gli identificativi di prenotazione e processo.

Questa combinazione deriva dai diversi requisiti delle operazioni. La
consultazione di una stazione o la validazione di un QR code richiedono un
esito immediato e sono quindi modellate come richieste REST. La telemetria,
invece, viene prodotta con continuità e non deve rallentare il simulatore:
viene perciò inviata su Kafka e consumata successivamente da Query. I
messaggi Camunda rappresentano eventi di dominio che fanno avanzare
un'istanza persistente del processo. L'identificativo della prenotazione
costituisce la chiave di correlazione e impedisce che un evento relativo a un
noleggio modifichi lo stato di un altro.

\capitoloNum{Il diagramma documentativo BPMN}

Durante questa fase abbiamo ripreso il sistema proiettato e l'abbiamo
riscritto come diagramma BPMN, aggiungendo tutti i calcoli e i controlli dei
vari servizi che non potevano essere espressi nella coreografia. Alcune
comunicazioni sono state inoltre modificate per rendere più semplice
l'implementazione. Il BPMN documentativo è suddiviso in cinque pool, una per
ogni ruolo individuato nell'analisi.

I due modelli sono contenuti nei file:
\begin{center}
\small
\path{docs/orchestration/modello_documentativo.bpmn}\\
\path{docs/orchestration/processes/process.bpmn}
\end{center}
Il primo descrive le responsabilità organizzative; il secondo traduce la pool
ACME in un processo Camunda 8, sostituendo le interazioni astratte con job
worker, connettori e messaggi correlati.

\section{Applicazione}

La pool App rimane in attesa di un messaggio da parte di ACME per ottenere la
lista delle stazioni e dei veicoli. Per gestire le due scelte dell'utente
abbiamo applicato un gateway esclusivo: visualizzare la mappa e avviare un
ritiro tramite QR code oppure prenotare un veicolo.

Nel ramo della mappa l'utente visualizza i veicoli e ne sceglie uno; queste
azioni sono rappresentate mediante user task. Successivamente invia ad ACME
il QR code acquisito. Nel ramo della prenotazione, dopo aver inviato la
richiesta, l'utente si trova di fronte a tre possibilità:
\begin{itemize}
  \item continuare con la prenotazione e notificare ACME al momento del ritiro;
  \item non ritirare il veicolo, lasciando che il timer di tolleranza scada;
  \item inviare una richiesta di annullamento.
\end{itemize}
I due rami che conducono al ritiro si ricongiungono nell'attesa della conferma
di ACME. Terminato l'uso del veicolo, l'applicazione invia la richiesta di
riconsegna e attende l'esito. Se la destinazione non è ancora stata raggiunta,
l'esito \code{VEHICLE\_NOT\_AT\_DESTINATION} mantiene il processo attivo e
consente un nuovo tentativo.

\section{Servizio stazioni}

La pool Stazione rimane in attesa di una notifica tramite un gateway basato
su eventi. Le due scelte sono la richiesta di blocco o quella di sblocco di
un veicolo. In entrambi i casi l'azione fisica è rappresentata nel modello
documentativo da una manual task; al termine viene notificato ad ACME l'esito
dell'operazione.

Nell'implementazione il servizio mantiene una sessione SOAP Jolie per ciascuna
stazione e offre inoltre un catalogo REST. In questo modo il comando sul
blocco viene separato dalla consultazione delle informazioni statiche.

\section{Servizio Fleet}

La pool Fleet rimane in attesa di un messaggio di avvio da parte di ACME, al
quale risponde notificandone la ricezione. Successivamente entra in un ciclo
di richieste dello stato del veicolo: ricerca e lettura dei dati sono
rappresentate mediante service task e l'esito viene inviato ad ACME. Quando
arriva la richiesta di arresto, Fleet interrompe la simulazione e restituisce
il riepilogo del viaggio.

Nel BPMN eseguibile un sottoprocesso parallelo effettua il monitoraggio ogni
due secondi e un evento di confine ne interrompe l'esecuzione alla
riconsegna. Posizione e batteria sono letture eventualmente consistenti,
mentre tempo e distanza calcolati dal simulatore sono i valori autorevoli per
l'addebito.

\section{Servizio bancario}

La pool Banca rimane in attesa di un messaggio tramite un gateway basato su
eventi. Il messaggio può riguardare:
\begin{itemize}
  \item un addebito, usato in caso di cancellazione tardiva o mancato ritiro;
  \item una cauzione, bloccata all'inizio del noleggio;
  \item il rilascio della cauzione, eseguito al termine del pagamento.
\end{itemize}
Ogni interazione con il conto è rappresentata mediante una service task e il
risultato viene notificato ad ACME. Se il saldo non è sufficiente,
l'applicazione richiede una ricarica e correla il relativo messaggio al
processo in attesa. Il servizio è volutamente dimostrativo e non rappresenta
un sistema di pagamento conforme ai requisiti di un sistema bancario reale.

\section{Servizio ACMEMobility}

La pool ACME era stata inizialmente suddivisa in due lane, per differenziare
la parte gestionale dalla parte logistica. La prima coordina il processo di
noleggio; la seconda gestisce stato, disponibilità e batteria dei veicoli.
In fase di implementazione tali responsabilità sono state separate nei
servizi ACMEMobility, Reservation, Logistics e Fleet.

La lane gestionale inizia inviando stazioni e veicoli all'applicazione, quindi
rimane in attesa della scelta fra QR code e prenotazione. Nel secondo caso un
gateway esclusivo gestisce tre continuazioni: ritiro, annullamento o scadenza
del tempo disponibile. I controlli temporali sono realizzati mediante timer
event. Quando avviene il ritiro, ACME richiede lo sblocco alla stazione,
blocca la cauzione e avvia Fleet, attendendo l'esito di ogni servizio.

Durante l'uso del veicolo viene richiesto periodicamente lo stato a Fleet. La
richiesta di riconsegna interrompe tale ciclo; ACME verifica la destinazione,
arresta Fleet, blocca il veicolo, calcola ed esegue il pagamento, rilascia la
cauzione e notifica l'applicazione.

Nell'implementazione ACMEMobility costituisce l'API pubblica e il confine di
sicurezza. Valida il JWT, verifica la proprietà delle prenotazioni, inoltra le
richieste ai servizi interni e pubblica i messaggi Camunda. Il processo
gestisce:
\begin{itemize}
  \item creazione, annullamento e scadenza della prenotazione;
  \item cauzione, ricarica eventuale, addebito e rilascio;
  \item sblocco, avvio Fleet, monitoraggio, blocco e arresto;
  \item calcolo del prezzo e notifiche SSE verso app e visualizzatore.
\end{itemize}
Timer e message boundary event modellano esplicitamente finestre temporali,
attese e interruzioni, evitando che tali regole restino nel frontend.

Il processo eseguibile realizza quindi una saga orchestrata. Non esiste una
transazione distribuita che comprenda Reservation, Logistics, Banca,
Stazione e Fleet: ogni partecipante completa una propria operazione locale,
mentre Camunda ne registra l'ordine e decide il passo successivo. Gli esiti
negativi restano espliciti nel flusso. Per esempio, un saldo insufficiente
porta all'attesa della ricarica e una riconsegna lontana dalla destinazione
non chiude la prenotazione. Questa soluzione rende visibile lo stato
intermedio del noleggio e consente di riprendere l'orchestrazione dopo un
errore temporaneo.

\section{Corrispondenza fra modello e implementazione}

La Tabella~\ref{tab:bpmn-implementation} riassume come gli elementi del
diagramma documentativo sono stati tradotti nel processo eseguibile. La
traduzione non è puramente sintattica: alcune attività astratte sono state
scomposte per rendere espliciti correlazione, attese e gestione degli esiti.

\begin{table}[ht]
\centering
\small
\begin{tabularx}{\textwidth}{@{}l X X@{}}
\toprule
Elemento & Implementazione Camunda & Motivazione \\
\midrule
Banca e Stazione & Connettore SOAP &
Riutilizzo delle interfacce WSDL dei servizi Jolie. \\
Connessione App--ACME & API REST e messaggi Camunda &
Richieste sincrone all'API e avanzamento asincrono del processo. \\
Calcolo del costo & Worker \code{calculate-charge} &
Logica Java testabile, separata dalle espressioni BPMN. \\
Monitoraggio Fleet & Sottoprocesso con timer di due secondi &
Aggiornamenti periodici senza bloccare il flusso principale. \\
Riconsegna & Messaggio correlato e boundary event &
Interruzione del monitoraggio della stessa prenotazione. \\
Attesa ricarica & Message catch event &
Sospensione persistente senza occupare un thread. \\
Scadenza & Timer event &
Applicazione deterministica delle finestre temporali. \\
Notifica & REST verso ACMEMobility, poi SSE &
Disaccoppiamento fra processo e connessione del browser. \\
\bottomrule
\end{tabularx}
\caption{Corrispondenza fra BPMN documentativo e processo eseguibile.}
\label{tab:bpmn-implementation}
\end{table}

\capitoloNum{Il diagramma SOA con UML}

\section{Decomposizione dei servizi}

In seguito abbiamo definito la struttura del sistema in un diagramma SOA
tramite UML, utilizzando il profilo TinySOA. Ogni servizio è stato definito
mediante un package, al quale è stato assegnato uno dei seguenti profili:
\begin{itemize}
  \item \emph{Task}, se il servizio implementa la logica del sistema;
  \item \emph{Entity}, se il servizio gestisce un'entità del dominio;
  \item \emph{Utility}, se il servizio è esterno alla logica applicativa.
\end{itemize}
Ogni package viene implementato da un'interfaccia con profilo
\emph{ServiceInterface} e da una classe con profilo \emph{Capability}. I due
elementi sono collegati da una dipendenza \emph{Expose}, diretta
dall'interfaccia alla classe. I servizi sono poi collegati mediante
dipendenze che partono dalla capability del consumatore e terminano
all'interfaccia del fornitore.

Oltre ai servizi già presenti nel diagramma BPMN, per semplificare
l'implementazione è stato necessario introdurre nuovi servizi oppure
suddividere un servizio in componenti più piccoli. In particolare sono stati
aggiunti il simulatore dei veicoli, utile a gestirne il movimento, il servizio
di autenticazione, il servizio Reservation e i componenti dedicati alla
telemetria. La decomposizione finale è riportata in
Tabella~\ref{tab:servizi}.

\begin{table}[ht]
\centering
\small
\begin{tabularx}{\textwidth}{@{}l l X@{}}
\toprule
Componente & Tipo & Responsabilità principale \\
\midrule
ACMEMobility & Task & API pubblica, autorizzazione, correlazione Camunda e SSE \\
Auth & Entity & utenti, credenziali, recupero password e JWT \\
Reservation & Entity & ciclo di vita e proprietà delle prenotazioni \\
Logistics & Entity & veicoli, disponibilità, QR e metadati batteria \\
Stations & Entity & catalogo stazioni e sessioni di blocco/sblocco \\
Banking & Task & saldo dimostrativo, cauzione, addebito e ricarica \\
Fleet & Task & coordinamento del viaggio e verifica della destinazione \\
Simulator & Task & percorso, movimento, durata, distanza e batteria finale \\
Tracking/Battery & Task & pubblicazione della telemetria su Kafka \\
Query & Entity & read model di posizione e batteria \\
Charge worker & Task & calcolo dell'importo corrente e finale \\
Camunda & Utility & stato e orchestrazione del processo BPMN \\
Kafka & Utility & trasporto asincrono della telemetria \\
GraphHopper & Utility & calcolo dei percorsi sulla mappa di Bologna \\
\bottomrule
\end{tabularx}
\caption{Servizi e responsabilità.}
\label{tab:servizi}
\end{table}

\section{Dipendenze e criteri progettuali}

Lo schema definisce App come consumatore di ACMEMobility. Quest'ultimo dipende
da Auth, Banking, Reservation, Logistics e Fleet, oltre a Camunda per
l'orchestrazione. Reservation gestisce la proprietà e il ciclo di vita delle
prenotazioni, mentre Logistics gestisce i veicoli. Fleet è a sua volta
dipendente da Query, Simulator, Reservation, Logistics e Stations. Simulator
usa GraphHopper e invia i dati ai servizi Tracking e Battery; questi
pubblicano su Kafka e Query ne costruisce il modello di lettura.

La separazione tra Reservation e Logistics impedisce che la disponibilità
del veicolo coincida con la proprietà della prenotazione. La scomposizione
Fleet--Simulator isola la simulazione dalla coordinazione. Tracking e Battery
scrivono eventi sui topic \code{vehicle-tracking} e
\code{vehicle-battery}; Query li combina in un read model, applicando un
semplice schema CQRS. GraphHopper e Kafka sono modellati come Utility perché
non sono servizi ideati nel progetto, ma dipendenze necessarie alla
comprensione dell'architettura. Il charge worker è un job worker Camunda
indipendente e può evolvere senza modificare il diagramma delle API.

La proprietà dei dati segue gli stessi confini. Auth è responsabile degli
utenti; Reservation delle prenotazioni; Logistics dell'inventario dei
veicoli; Stations del catalogo delle stazioni. Nessun servizio modifica
direttamente le tabelle di un altro dominio: le variazioni passano attraverso
le interfacce esposte. Il read model di Query non sostituisce le sorgenti
autorevoli, ma conserva una proiezione ricostruibile degli eventi di
posizione e batteria. La scelta riduce l'accoppiamento sullo schema e
consente ai servizi di evolvere in modo indipendente.

\capitoloNum{Implementazione}

\section{Servizio bancario}

Il servizio bancario è stato implementato con Jolie, impostando una
\emph{inputPort} con protocollo SOAP. Nel server sono state definite le
seguenti operazioni:
\begin{itemize}
  \item \code{login} e \code{logout}, per creare e terminare l'ambiente
        privato dell'utente;
  \item \code{charge}, per sottrarre l'importo richiesto dal saldo;
  \item \code{holdDeposit}, per bloccare la cauzione e memorizzarla nella
        sessione privata;
  \item \code{releaseDeposit}, per restituire la cauzione;
  \item \code{checkFunds}, per verificare la disponibilità prima
        dell'addebito.
\end{itemize}

La ricarica autorizzata dall'utente non appartiene alla sessione SOAP:
ACMEMobility invoca l'operazione REST \code{rechargeAccount}, esposta dalla
stessa applicazione Jolie.

Il server rimane in attesa di un login e, dopo l'arrivo della richiesta,
crea una sessione personale per l'utente. All'interno della sessione rimane
in attesa delle operazioni sul conto fino alla ricezione del logout. Rispetto
al modello iniziale abbiamo corretto la gestione del saldo insufficiente: il
servizio non effettua autonomamente una ricarica senza autorizzazione, ma
restituisce l'esito al processo. ACMEMobility notifica l'applicazione e
Camunda attende un messaggio di ricarica correlato alla prenotazione.

\begin{sloppypar}
L'interfaccia è descritta nel file \code{BankingService.wsdl}, mentre
\code{banking-\allowbreak interface.iol} contiene i tipi e le operazioni
condivise. Il file \code{banking-\allowbreak service.ol} implementa il
comportamento conversazionale.
Nel deployment Kubernetes, connessione al database e parametri del servizio
sono passati al container tramite configurazione. Il BPMN usa invece il nome
DNS stabile \code{banking-service.acme} per i connettori interni. Il servizio
rimane dimostrativo: saldo e credenziali
non hanno le garanzie di sicurezza di un vero intermediario finanziario.
\end{sloppypar}

\section{Servizio stazioni}

Il servizio stazioni è stato implementato con Jolie, come il servizio
bancario. Utilizza due interfacce differenti: una SOAP per comunicare con il
processo Camunda e una REST per fornire a Logistics il catalogo delle
stazioni. Nel server sono state definite le funzioni:
\begin{itemize}
  \item \code{openSession} e \code{closeSession}, per creare e terminare le
        sessioni delle singole stazioni;
  \item \code{unlockVehicle} e \code{lockVehicle}, per rappresentare lo
        sblocco e il blocco di un veicolo;
  \item \code{listStations}, che restituisce l'elenco delle stazioni;
  \item \code{getStation}, che restituisce la stazione richiesta.
\end{itemize}

Il server rimane in attesa dell'apertura di una sessione, della richiesta di
una stazione o della richiesta dell'intera lista. Quando arriva una
\code{openSession} viene aperta una sessione privata, nella quale possono
essere ricevuti più comandi di blocco e sblocco. La \code{closeSession}
termina la sessione e riporta il servizio nello stato iniziale. Il catalogo e
le sessioni delle stazioni, incluso
l'identificativo dell'ultimo veicolo gestito, sono persistiti in PostgreSQL
tramite il driver JDBC di Jolie; questo evita che la perdita del container
comporti la perdita dei dati.

La separazione tra SOAP e REST riflette due tipi di interazione differenti.
Il comando di blocco appartiene a una conversazione con stato, mentre la
lettura del catalogo è una richiesta priva di stato e può essere esposta come
risorsa. In questo modo Logistics non dipende dalla sessione usata da
Camunda.

\section{Servizio Fleet management}

Come anticipato nel diagramma UML, il servizio Fleet management è suddiviso
nei componenti Fleet, Simulator, Tracking, Battery e Query. Tutti questi
componenti sono implementati in Java 21 con Helidon MicroProfile 4.3.4.

Il \code{tracking-service} riceve dal simulatore la posizione del veicolo e
la pubblica sul topic Kafka \code{vehicle-tracking}. In modo analogo, il
\code{battery-service} pubblica la percentuale di carica sul topic
\code{vehicle-battery}. I producer richiedono conferma da tutte le repliche
con \code{acks=all}, abilitano l'idempotenza e configurano retry e timeout.
Queste impostazioni riducono il rischio di perdita o duplicazione degli
eventi durante errori transitori.

Il \code{query-service} consuma entrambi i topic e mantiene l'ultima
fotografia disponibile per ciascun veicolo. In questo modo viene realizzato
un semplice modello CQRS: Tracking e Battery costituiscono il lato di
scrittura degli eventi, mentre Query costituisce il lato di lettura. Le due
misure possono arrivare in istanti differenti; di conseguenza posizione e
batteria mostrate nell'interfaccia sono eventualmente consistenti.

Ogni evento contiene l'identificativo del veicolo e il valore misurato. Query
usa tale identificativo come chiave logica e aggiorna soltanto la parte del
modello interessata dall'evento ricevuto. Un aggiornamento di posizione non
sovrascrive quindi la batteria e viceversa. La configurazione idempotente dei
producer riduce i duplicati dovuti ai retry; il consumatore, memorizzando
l'ultima fotografia, rende comunque innocua la ripetizione dello stesso
valore. Questa consistenza eventuale è accettabile per la visualizzazione,
ma non per il calcolo economico: minuti, chilometri e batteria finale vengono
richiesti direttamente al simulatore al termine del viaggio.

Il \code{fleet-service} coordina il sottosistema. Per iniziare un viaggio
recupera i dati della prenotazione e del veicolo, seleziona la destinazione e
richiede al simulatore di iniziare il percorso. Durante il viaggio interroga
Query e Simulator per produrre lo stato richiesto dal BPMN. Alla
riconsegna controlla che il veicolo abbia raggiunto la stazione assegnata:
in caso contrario restituisce l'esito
\code{VEHICLE\_NOT\_AT\_DESTINATION}; in caso positivo arresta la simulazione
e restituisce minuti, chilometri e batteria finale.

\section{Servizio ACMEMobility}

La parte gestionale di ACME è stata realizzata come BPMS tramite Camunda 8.
Per farlo abbiamo ripreso il BPMN documentativo, isolando la parte gestionale
e trasformandola in un processo eseguibile. Le coppie di task che comunicano
con i servizi Jolie sono state sostituite da connettori SOAP; le
comunicazioni con API HTTP sono state sostituite da connettori REST o da job
worker. Le variabili del processo contengono identificativi, tempi e risultati
necessari alla correlazione. Il BPMN inizializza gli endpoint con i nomi DNS
stabili dei Service Kubernetes; credenziali e segreti restano invece nella
configurazione del deployment.

Il ciclo di lettura dello stato del veicolo è stato spostato in un
sottoprocesso parallelo, così da poterlo gestire indipendentemente dal flusso
principale. Il ciclo viene attivato da un timer ogni due secondi e termina
mediante un evento di confine quando arriva la richiesta di riconsegna.
Questa soluzione sostituisce il gateway basato su eventi del diagramma
documentativo e permette al processo di continuare a ricevere messaggi mentre
il monitoraggio è in corso.

ACMEMobility è inoltre un microservizio Helidon che funge da API pubblica.
Espone le operazioni di autenticazione, stazioni, prenotazioni, QR, banca e
inventario dei veicoli; valida il JWT, ricava l'identità dell'utente e
impedisce l'accesso alle prenotazioni altrui. Fleet è invece invocato dal
processo e dagli altri servizi interni. Il client Camunda avvia istanze e
pubblica messaggi correlati.

Per notificare il browser sono stati usati Server-Sent Events. Sono presenti
flussi distinti per prenotazioni, veicoli e visualizzatore della simulazione.
La comunicazione è principalmente server-to-client e non richiede la
bidirezionalità di WebSocket; SSE permette inoltre al browser di ristabilire
la connessione in modo semplice.

Il calcolo dell'addebito è affidato al job worker
\code{calculate-charge}, scritto in Java. Esso riceve minuti, chilometri e
batteria e applica la formula
\[
  C = 1{,}00 + 0{,}20\,m + 0{,}50\,k,
\]
dove $m$ rappresenta i minuti e $k$ i chilometri. Se la batteria è inferiore
al 15\%, il costo viene moltiplicato per $1{,}10$; l'importo finale è
arrotondato ai centesimi. Il worker valida inoltre che minuti e chilometri
siano non negativi e che la batteria appartenga all'intervallo $[0,100]$.
Un secondo worker, \code{stop-and-charge}, gestisce la ricarica di sicurezza:
quando la batteria raggiunge la soglia critica del 5\%, richiede a Fleet di
ricaricare il veicolo al 25\% e di riprendere la simulazione ancora attiva.

La separazione del calcolo in un worker evita di inserire codice complesso
nelle espressioni BPMN e rende la regola tariffaria verificabile tramite test
unitari. Il worker restituisce sia l'importo base sia quello finale e indica
esplicitamente se la maggiorazione è stata applicata. Il processo può così
notificare all'utente il motivo della differenza, mantenendo distinta la
decisione economica dall'operazione SOAP che modifica il saldo.

\section{Servizio di autenticazione}

Il servizio di autenticazione è implementato con Helidon e separa la gestione
degli utenti dal resto del dominio. Espone registrazione, login, cambio
password e recupero password. Le credenziali e i dati dell'utente sono
persistiti in PostgreSQL; Flyway crea e aggiorna lo schema in fase di avvio.
Per i test è disponibile un repository in memoria, così da non rendere le
verifiche unitarie dipendenti dal database.

Dopo un login valido il servizio genera un JWT. Il token viene inviato dal
frontend nell'header \code{Authorization} e verificato da ACMEMobility prima
di inoltrare una richiesta. In questo modo i servizi interni non sono esposti
direttamente al browser. Il recupero password usa un'astrazione
\code{PasswordMailer}, con implementazione SMTP e sostituto controllabile nei
test. Le configurazioni incluse nel repository sono adatte alla
dimostrazione, ma chiavi, password e segreti dovrebbero essere ruotati e
gestiti esternamente in un ambiente produttivo.

\section{Servizio Reservation}

Reservation gestisce la proprietà e il ciclo di vita della prenotazione. I
dati principali sono identificativo, utente, veicolo, stazioni, orario
previsto e stato. La persistenza usa PostgreSQL, pool HikariCP e migrazioni
Flyway. Le transizioni sono esposte mediante API REST interne, invocate dai
connettori Camunda, in modo che il processo possa creare, ritirare, annullare,
far scadere e completare una prenotazione.

La prenotazione può iniziare nell'intervallo compreso tra l'istante corrente e
i 30 minuti successivi. L'annullamento è gratuito fino a cinque minuti prima
dell'orario stabilito; dopo tale soglia viene applicata la penale di 10 euro.
Anche il mancato ritiro entro i cinque minuti di tolleranza produce lo stesso
addebito. Il servizio associa sempre la prenotazione all'utente autenticato;
ACMEMobility verifica tale proprietà prima di pubblicare comandi Camunda.

\section{Servizio Logistics}

Logistics gestisce inventario e stato operativo dei veicoli: stazione,
disponibilità, batteria, coordinate e QR code. I dati sono memorizzati in
PostgreSQL e inizializzati tramite migrazioni Flyway. Il servizio consulta il
catalogo REST di Stations per associare correttamente i mezzi alle stazioni.

Quando il processo conferma una prenotazione, Logistics pone il veicolo in
stato riservato; al ritiro lo rende in uso; al termine aggiorna stazione,
batteria e disponibilità. La separazione da Reservation evita che i dati relativi alla
proprietà del noleggio vengano confusi con quelli dell'inventario. Nel flusso
QR, ACMEMobility confronta il codice inserito con il valore restituito da
Logistics prima di avviare il processo.

\section{Simulazione dei veicoli}

Il simulatore è un servizio Helidon che riproduce lo spostamento del veicolo.
All'avvio riceve posizione iniziale e destinazione, quindi usa GraphHopper e
la mappa di Bologna inclusa nel progetto per calcolare un percorso stradale.
Il movimento viene avanzato per passi; a ogni aggiornamento il simulatore
invia la posizione a Tracking e la carica a Battery.

Il simulatore mantiene i valori autorevoli di tempo trascorso e distanza
percorsa. Fleet li usa sia durante il monitoraggio sia nel riepilogo finale,
evitando che il costo dipenda dal numero di richieste del frontend. Quando la
batteria raggiunge una soglia critica il processo può fermare il veicolo,
richiedere una ricarica e riprendere la simulazione. Il visualizzatore mostra
separatamente percorso, posizione, batteria, durata, chilometri e importo
corrente.

\section{Applicazione}

L'applicazione è stata sviluppata come single-page application mediante
React, TypeScript, HTML e CSS, con Vite come strumento di build. Le viste
gestiscono login, registrazione, mappa, elenco delle prenotazioni, ritiro,
annullamento e riconsegna. L'utente può selezionare una stazione e un veicolo,
indicare l'orario di inizio entro la finestra ammessa oppure inserire il QR
code per il ritiro immediato.

Le notifiche SSE aggiornano l'interfaccia senza polling continuo e mostrano
conferma, cauzione, eventuale penale, stato del viaggio e pagamento finale.
Una seconda applicazione React, il simulation visualizer, combina la
posizione e la batteria lette da Query con durata, chilometri e costo inviati
da ACMEMobility. La distinzione evita di accoppiare la UI operativa a quella
di osservazione.

\section{Container e Kubernetes}

Ogni servizio dispone di un Dockerfile. Gli script nella directory
\code{scripts} preparano un profilo Minikube, creano i namespace
\code{acme}, \code{orchestration} e \code{kafka}, installano Camunda tramite
Helm e Kafka tramite Strimzi, costruiscono le immagini all'interno di
Minikube e attendono la readiness dei deployment.

Kafka usa i topic \code{vehicle-tracking} e \code{vehicle-battery};
PostgreSQL è distribuito nel namespace applicativo; GraphHopper usa la mappa
inclusa nel repository. L'API ACMEMobility e il visualizzatore sono esposti
sulle NodePort fisse 30080 e 30081. Anche Query è attualmente dichiarato come
\code{NodePort}, con porta esterna assegnata dal cluster, ma non viene
contattato direttamente dai due frontend. Secret ed environment variable
separano credenziali e configurazione dalle immagini.

I servizi applicativi sono raggiunti mediante nomi DNS interni stabili,
indipendenti dall'indirizzo dei singoli pod. Per il normale utilizzo, il
browser comunica soltanto con ACMEMobility e con il visualizzatore; database
e broker non vengono pubblicati all'esterno. Gli script attendono inoltre il rollout
dei deployment prima di avviare le verifiche, evitando che una dipendenza
ancora in fase di inizializzazione venga interpretata come errore
applicativo.

La procedura completa richiede risorse elevate soprattutto a causa di
Camunda, Elasticsearch, Kafka e GraphHopper. La configurazione proposta
assegna 16 GiB di memoria e 8 CPU al profilo Minikube. Gli script non
eliminano automaticamente un profilo esistente con risorse insufficienti,
per evitare la perdita accidentale dello stato dell'utente.

\section{Gestione dei fallimenti}

L'orchestrazione riduce il rischio che un errore temporaneo renda
incomprensibile lo stato del noleggio, ma non equivale a una transazione
distribuita. Se un servizio non risponde, l'istanza Camunda conserva le
variabili e il punto di avanzamento; il job o il connettore può essere
ritentato senza ricostruire il flusso dal browser. I comandi sono correlati
alla prenotazione e le verifiche di stato nei servizi limitano l'effetto di
messaggi ripetuti.

Rimangono casi che, in un sistema produttivo, richiederebbero compensazioni
più articolate. Se la cauzione viene bloccata ma lo sblocco fallisce
definitivamente, il processo dovrebbe rilasciarla esplicitamente. Se Fleet si
arresta dopo l'avvio, occorre ricostruire lo stato dal simulatore prima di
accettare la riconsegna. Un ritardo di Kafka può rendere temporaneamente
obsoleta la posizione mostrata, ma non modifica l'addebito, che usa il
riepilogo autorevole del simulatore. Per completare l'industrializzazione
sarebbero inoltre utili dead-letter queue, politiche di retry differenziate,
timeout per ogni dipendenza e operazioni amministrative per correggere
un'istanza bloccata.

\section{Verifica}

La verifica combina test automatici e scenari end-to-end. I test JUnit
coprono le risorse Helidon, le transizioni principali e il calcolo
dell'addebito. Per i frontend la verifica automatica comprende la compilazione
TypeScript, la build di produzione e l'analisi statica con ESLint; non sono
presenti test automatici dei componenti React. Le procedure manuali di smoke
test descritte nel README verificano autenticazione, prenotazione, ritiro,
telemetria, riconsegna, cancellazione e ricarica.

Le prove sono organizzate su livelli differenti. I test unitari isolano le
regole deterministiche, come formula tariffaria, validazione degli input e
comportamento dei repository sostitutivi. I test dei singoli servizi
verificano contratti HTTP, codici di risposta e serializzazione. Le procedure
di smoke test vengono invece eseguite manualmente sul cluster completo e
controllano che Camunda, Jolie, PostgreSQL, Kafka e i servizi Helidon
collaborino nello stesso scenario. La prova manuale del frontend completa la
verifica osservando notifiche SSE, aggiornamento della mappa e transizioni
visibili all'utente.

Questa distinzione è importante perché un test unitario non può dimostrare la
corretta correlazione di un messaggio Camunda o la consegna della telemetria
attraverso Kafka. Simmetricamente, uno scenario end-to-end è più lento e
rende meno immediata l'individuazione di un errore nella formula. L'uso
congiunto dei livelli bilancia rapidità diagnostica e copertura
dell'integrazione.

I criteri di accettazione principali sono:
\begin{itemize}
  \item una prenotazione riserva un solo veicolo e rispetta le finestre
        temporali;
  \item un QR errato non avvia il processo e un utente non può operare sulle
        prenotazioni altrui;
  \item posizione e batteria raggiungono Query attraverso Kafka;
  \item una richiesta di riconsegna anticipata non chiude il viaggio;
  \item addebito, maggiorazione e cauzione concordano con il riepilogo finale;
  \item tutti i deployment applicativi raggiungono lo stato \emph{Ready}.
\end{itemize}

\capitoloNum{Conclusioni}

Il progetto realizza l'intero ciclo di un noleggio mediante un'architettura
eterogenea a microservizi. L'orchestrazione BPMN rende esplicite regole,
timeout e compensazioni; REST e SOAP gestiscono le richieste con risposta,
Kafka disaccoppia la telemetria e SSE aggiorna le interfacce. PostgreSQL
assegna la proprietà dei dati ai rispettivi servizi, mentre Kubernetes rende
riproducibile il deployment.

Rispetto a un'applicazione monolitica, questa soluzione introduce un costo
operativo evidente: più processi, configurazioni, immagini, protocolli e
dipendenze devono essere distribuiti e osservati. Una singola operazione di
noleggio attraversa più confini di rete e non può affidarsi a una sola
transazione ACID. In compenso, i confini rendono esplicita la proprietà dei
dati, consentono di sostituire o scalare separatamente i componenti e
permettono di integrare tecnologie eterogenee, come Jolie, Camunda e Helidon,
attraverso contratti definiti.

Nel contesto del progetto la decomposizione è volutamente più articolata di
quanto richiederebbe il solo carico dimostrativo. Il suo valore risiede nella
possibilità di studiare orchestrazione, consistenza eventuale,
interoperabilità e deployment cloud-native su un caso d'uso completo. Per
un prodotto con scala ridotta, alcuni servizi potrebbero essere inizialmente
accorpati e separati soltanto in presenza di esigenze organizzative o di
carico misurabili.

Il principale compromesso è la complessità operativa: per un prototipo, il
numero di componenti e il consumo di risorse sono superiori a quelli di
un'applicazione monolitica. Tale costo è giustificato dall'obiettivo
didattico di integrare orchestrazione, messaggistica, protocolli eterogenei e
deployment cloud-native. Restano fuori dall'ambito produttivo la gestione
sicura dei segreti, l'alta disponibilità, la cifratura end-to-end,
l'osservabilità centralizzata e un vero payment provider. Questi aspetti
costituiscono le evoluzioni naturali del sistema.

\end{document}
